\section{Brook Taylor und Joseph Fourier}

\begin{ejercicio}[Konvergenzradius einer Potenzreihe]
    Berechnen Sie den Konvergenzradius $r$ der Potenzreihe
    \begin{equation*}
        s = \sum_{n=2}^{\infty} \dfrac{x^n}{\sqrt{n^2 - 1}}
    \end{equation*}
\end{ejercicio}

\begin{ejercicio}[Taylor-Entwicklung mit Wurzelfunktion]
    Berechnen Sie die Taylorsche Reihe der Funktion
    \begin{equation*}
        f(x) = \dfrac{1}{\sqrt{x}}
    \end{equation*}
    am Entwicklungspunkt $x_0 = 1$ bis zur 4. Ordnung und bestimmen Sie den Konvergenzbereich.
\end{ejercicio}

\begin{ejercicio}[Relativistische Massenzunahme]
    Die Masse $m$ eines Elektrons nimmt nach der Relativitätstheorie mit der Geschwindigkeit $v$ zu. Es gilt:
    \begin{equation*}
        m = m(v) = \dfrac{m_0}{\sqrt{1 - \left(\dfrac{v}{c}\right)^2}}
    \end{equation*}
    wobei $m_0$ die Ruhemasse des Elektrons und $c$ die Lichtgeschwindigkeit ist. Entwickeln Sie mit Hilfe der Potenzreihenentwicklung eine Näherungsformel für die Abhängigkeit zwischen Masse und Geschwindigkeit unter der Annahme $v \ll c$.
\end{ejercicio}

\begin{ejercicio}[Freier Fall mit Reibung]
    Unter Berücksichtigung des Luftwiderstandes besteht zwischen der Fallgeschwindigkeit $v$ und dem Fallweg $s$ der folgende (komplizierte) Zusammenhang:
    \begin{equation*}
        v = v(s) = \sqrt{\dfrac{mg}{k} \left(1 - e^{-\frac{2ks}{m}}\right)}
    \end{equation*}
    wobei $m$ die Masse des Körpers, $g$ die Erdbeschleunigung und $k$ der Reibungskoeffizient (mit $k > 0$) ist. Wie lautet dieses Fallgesetz im luftleeren Raum?
\end{ejercicio}

\begin{ejercicio}[``Angeschnittener'' Wechselstrom]
    Ein ``angeschnittener'' Wechselstrom, dessen Zeitabhängigkeit im Periodenintervall $0 \leq t \leq T$ durch die Gleichung
    \begin{equation*}
        i(t) =
        \begin{cases}
            \hat{i} \cdot \cos (\omega_0 t) & 0 \leq t \leq \dfrac{T}{4} \\
            0 & \dfrac{T}{4} \leq t \leq T
        \end{cases}
    \end{equation*}
    beschrieben wird (Kreisfrequenz $\omega_0 = \dfrac{2\pi}{T}$). Wie lautet die Fourier-Zerlegung dieser Funktion in reeller Darstellung?
\end{ejercicio}